\documentclass[letterpaper,10pt,conference]{ieeeconf}

\IEEEoverridecommandlockouts
\overrideIEEEmargins
\usepackage{amsmath}
\usepackage{graphicx}
\usepackage{svg}

\title{PIRoute: Supplementary Figures}
\author{Anonymous Authors}

\begin{document}
\maketitle
\thispagestyle{empty}
\pagestyle{empty}

This document contains the supplementary figures for the PIRoute manuscript.
The figures are exploratory or qualitative and are not pooled with the G25
confirmatory analysis.

\begin{figure*}[t]
  \centering
  \includesvg[inkscapelatex=false,width=0.92\textwidth]{generated/figS1_generalization/generalization}
  \caption{Zero-update robot-count check. PIRoute is transferred without actor,
  router, or threshold updates to three- and seven-robot scenes. The full-team
  success differences are exploratory and their confidence intervals cross zero.}
  \label{fig:s1}
\end{figure*}

\begin{figure*}[t]
  \centering
  \includesvg[inkscapelatex=false,width=0.92\textwidth]{generated/figS2_trajectories/trajectory_overview}
  \caption{Representative closed-loop trajectories from the post-sealed
  qualitative capture. The outcome-stratified sample is used for visual
  diagnosis and not for re-estimating episode-level rates.}
  \label{fig:s2}
\end{figure*}

\begin{figure*}[t]
  \centering
  \includesvg[inkscapelatex=false,width=0.92\textwidth]{generated/figS3_gate_timeline/gate_timeline}
  \caption{Router probability and executed-policy timelines for the same
  qualitative capture scenes as Fig.~\ref{fig:s2}. The Router is updated every
  two environment steps; the timeline is descriptive rather than a population
  estimate of switching frequency.}
  \label{fig:s3}
\end{figure*}

\begin{figure*}[t]
  \centering
  \includesvg[inkscapelatex=false,width=0.92\textwidth]{generated/figS4_external_router/g26_e1_effects}
  \caption{Exploratory comparison with a normalizing-flow-inspired external
  switching baseline. This baseline is a local implementation inspired by prior
  work, not a strict reproduction, and is not pooled with G25.}
  \label{fig:s4}
\end{figure*}

\end{document}
